\chapter{System Study}

The current landscape of online security tools offers limited, fragmented support for everyday users who need to quickly assess the trustworthiness of a website. Most existing solutions either rely on a single signal (such as a blacklist check or SSL badge) or require technical knowledge to interpret. Through research and observation, we identified the need for a unified, multi-signal platform that is accessible to non-technical users while being thorough enough to catch sophisticated threats. The ProblySUS system study outlines the limitations of existing tools and how our proposed system addresses them.

\section{Existing System}

Most current URL safety tools — such as Google Safe Browsing, VirusTotal URL scanner, or built-in browser warnings — provide basic threat detection. However, these platforms:

\begin{itemize}
\item Rely primarily on blacklist lookups and miss zero-day phishing sites not yet reported.
\item Do not perform live behavioral analysis of how a page actually behaves in a browser.
\item Provide no insight into domain age, DNS infrastructure gaps, or content-level trust indicators.
\item Offer no privacy or tracker auditing — users receive no information about how aggressively a site profiles them.
\item Present results as simple pass/fail flags with no explainable reasoning, making it difficult for users to understand why a site is flagged.
\end{itemize}

These systems function as reactive lookup tools rather than proactive, intelligent assessors of website trustworthiness.

\section{Proposed System}

ProblySUS is an end-to-end, multi-signal scam and phishing detection platform that analyzes any submitted URL through ten independent analysis modules and produces a composite risk score (0–100) with a label — \textit{Safe}, \textit{Caution}, \textit{Suspicious}, or \textit{Fraudulent} — backed by a human-readable list of reasons.

\textbf{Core Advantages:}

\begin{itemize}
\item Multi-signal analysis combining domain intelligence, URL patterns, DNS records, live content, blacklist lookup, and headless browser simulation.
\item Runtime behavioral analysis via Playwright that detects JavaScript-driven redirects and dynamic external requests invisible to static scanners.
\item Privacy and tracker auditing revealing cookie tracking intensity, fingerprinting techniques, and named commercial trackers.
\item Fully explainable output — every score point is justified with a plain-language reason.
\item Accessible through both a web application and a Chrome browser extension for seamless in-browser scanning.
\item Local-first architecture requiring no third-party API keys, with an auto-refreshing URLhaus threat feed.
\end{itemize}

These functionalities are modular, enabling independent maintenance and future extension of any analysis component without affecting the rest of the pipeline.

\section{Software Requirements Specification}

ProblySUS is designed to empower everyday internet users with an intelligent, transparent tool for assessing website safety. The system combines heuristic, behavioral, and content-based signals into a unified risk assessment. The application is intended for a broad audience including students, professionals, and non-technical users who frequently encounter unfamiliar URLs.

The system accepts a URL as input and generates a comprehensive risk report including a numeric score, risk label, ordered list of reasons, SSL/HTTPS status, domain age, DNS infrastructure details, blacklist status, browser behavior summary, network domain classification, privacy grade, and identified trackers. The platform is accessible via a React-based web application and a Chrome browser extension, both connected to the same Flask REST API backend.

\subsection{Overall Description}

ProblySUS operates as a self-contained platform with a decoupled frontend and backend communicating through a REST API. The backend runs locally on port 5000 while the frontend operates on port 5173, with CORS enabled for cross-origin communication. The Chrome extension connects directly to the local backend, allowing both interfaces to use the same analysis engine.

The system maintains a local URLhaus threat intelligence database which automatically refreshes every 24 hours via a background thread. This ensures that the threat database remains current without requiring manual updates. Future versions are expected to support cloud deployment for public access and integration with browser-level safe browsing APIs for real-time protection.

\subsection{Product Functions}

The system includes several core functionalities that together provide a comprehensive risk assessment for any submitted URL.

\textbf{URL Validation and Normalization}

Every submitted URL is validated for syntactic correctness. Missing schemes such as \texttt{http://} or \texttt{https://} are automatically added. The system also checks HTTPS usage and TLS certificate validity before continuing further analysis.

\textbf{Domain Intelligence}

WHOIS registry data is queried to determine the domain age in days. Domains registered within the past 30 days are flagged as high-risk signals. DNS records including MX and SPF are also checked, since the absence of email infrastructure often indicates non-legitimate domains.

\textbf{URL Pattern Analysis}

The URL structure is inspected for suspicious characteristics such as IP-based hostnames, abused top-level domains (e.g., \texttt{.tk}, \texttt{.xyz}, \texttt{.pw}), excessive hyphens, and phishing keywords.

\textbf{Content Trust Analysis}

The system fetches live page HTML and scans it for urgency phrases, sensitive data request forms, presence of trust pages (Privacy Policy, Contact, About), external link ratios, and reading complexity scores.

\textbf{Blacklist Lookup}

The submitted domain is compared against the locally cached URLhaus threat feed. If a match is found, the system immediately returns a \textit{Fraudulent} verdict with a score of 100.

\textbf{Behavioral Analysis}

A headless Chromium browser launched via Playwright loads the target page, executes JavaScript, and captures redirect chains and outgoing network requests. This allows the system to detect behavior that static scanners cannot observe.

\textbf{Network, Privacy, and Tracker Profiling}

External domains contacted during page execution are categorized into safe infrastructure, suspicious, or unknown groups. Cookie usage, third-party scripts, and browser fingerprinting techniques are analyzed. Known commercial trackers are identified using a local tracker database.

\textbf{Risk Scoring}

Outputs from all modules are processed by \texttt{scorer.py}, which applies a confidence-weighted penalty model and compound signal amplification to generate the final score and explanatory reasons.

\section{Requirements Specification}

This section describes the functional and non-functional requirements of ProblySUS including system interfaces, performance expectations, and architectural design.

\subsection{External Interfaces}

\textbf{User Interface}

The web application includes a full-page analysis dashboard with a URL input field, animated risk score gauge, color-coded label badge, expandable reasons list, and detailed panels for behavior, network activity, privacy analysis, and tracker detection.

The Chrome extension provides a lightweight popup interface displaying the score arc, key metrics (SSL, domain age, blacklist, MX), risk label, and tabbed views for detailed information.

\textbf{API Interface}

The frontend communicates with the Flask backend through REST APIs.

\begin{itemize}
\item \texttt{POST /analyze} — accepts a URL and returns the full analysis report in JSON format.
\item \texttt{POST /api/blacklist/update} — manually triggers an update of the URLhaus blacklist.
\end{itemize}

Optional parameters include \texttt{?fast=true} to skip Playwright analysis and \texttt{?force=true} to bypass cached results.

\textbf{Hardware Requirements}

The system requires a machine capable of running Python 3 and Node.js. The Playwright module requires approximately 170 MB for the Chromium browser binary and an active internet connection. A minimum of 4 GB RAM is recommended.

\textbf{Software Compatibility}

The web application supports modern browsers including Chrome, Firefox, Edge, and Safari. The extension supports Chromium-based browsers such as Chrome, Edge, and Brave.

The backend requires Python 3.8+ and dependencies including Flask, flask-cors, requests, python-whois, beautifulsoup4, tldextract, dnspython, textstat, and playwright.

\subsection{Functional Requirements}

\begin{itemize}
\item The system shall accept syntactically valid URLs and normalize missing schemes.
\item The system shall validate HTTPS usage and TLS certificate integrity.
\item The system shall retrieve WHOIS data and classify domain age into risk tiers.
\item The system shall check DNS records for MX and SPF entries.
\item The system shall perform static URL pattern analysis.
\item The system shall analyze live page HTML for trust and phishing indicators.
\item The system shall check the domain against the URLhaus blacklist.
\item The system shall execute JavaScript using a headless browser for behavioral analysis.
\item The system shall analyze external network requests and classify domains.
\item The system shall identify commercial trackers and fingerprinting techniques.
\item The system shall compute a final composite risk score and assign a risk label.
\item The system shall cache analysis results for 30 minutes.
\item The system shall update the URLhaus threat feed automatically every 24 hours.
\end{itemize}

\subsection{Non-Functional Requirements}

\textbf{Performance}

Standard URL analysis should complete within 5–6 seconds. Full analysis including browser simulation should complete within 8–12 seconds. Fast mode analysis should complete within 3–5 seconds. Cached results should be returned instantly.

\textbf{Usability}

The system is designed for non-technical users. Color-coded labels (green, yellow, orange, red), visual indicators, and plain-language explanations ensure that results are easy to interpret.

\textbf{Reliability}

If the Playwright module fails or is unavailable, the system gracefully degrades while continuing other analyses. All modules implement exception handling to prevent system crashes.

\textbf{Maintainability}

Each analysis function is implemented as an independent module following the single-responsibility principle. New signals can be integrated by adding a new module and extending the scoring logic.

\textbf{Scalability}

The system supports parallel execution of analysis modules using a \texttt{ThreadPoolExecutor}. Result caching reduces redundant processing. The architecture supports future deployment using containerization technologies such as Docker.

\textbf{Portability}

The system runs on Windows, macOS, and Linux without modification. The frontend is a React/Vite application deployable to any static hosting service, while the backend is a standard Flask application.